\section{Statement of Patient Commitment}

The parent clinical protocol opens with a Statement of Compliance addressed to
the Food and Drug Administration, the Institutional Review Board, and the
sponsor \cite{onpremwhippl}. That statement is necessary, it is correct, and it
is not addressed to the person whose abdomen will be opened. It tells the
regulator which parts of Title 21 the study satisfies. It does not tell the
participant what the study owes them.

This paper opens instead with the seven things the same protocol commits to the
participant. Each one is tied to a clause that makes it enforceable rather than
aspirational, and each one names something the participant can check. The
distinction matters, because the surveyed literature on patient attitudes to
robotic and artificial-intelligence-assisted surgery reports again and again
that reassurance without a checkable referent is heard as marketing
\cite{Jauniaux2025, BrarRAS2024X}.

The premise is taken from the author's earlier analysis of proposed United
States legislation for Physical AI oncology trials: the cancer patient, not the
doctor, not the nurse, not the trial sponsor, not the Institutional Review
Board, and not the regulator, is the priority participant
\cite{paibillprior, hr9510billv5}. That premise is not a sentiment. It has a
structural consequence, which the whole of this paper works out: if the
participant is the priority, then every safeguard written as a sponsor
obligation can also be written as a participant capability, and the two
readings should agree. Where they do not agree, this paper says so.

This is an advocacy document assembled from the protocol's own record. It is not
the consent form and it is not a substitute for one. It is designed to be read
before the consent conversation, so that the conversation can be spent on the
questions this document cannot answer.

\subsection{The seven commitments}

\textbf{One. Nothing starts without you.} No study procedure occurs before
informed consent is documented under 21 CFR \S 50.20 \cite{ecfr2026part50}.
That includes the staging imaging repeated for study purposes, the tissue
re-review, and every laboratory draw the protocol adds beyond routine care. The
check is simple: the consent form carries a date, and no study record may
predate it.

\textbf{Two. The robot is proven before you.} A Phase 0 simulation-validation
gate precedes any patient-facing robotic activity: at least 1000 simulated
procedures across at least two independent simulation frameworks, reproducing
motion within a 2 mm positional and 0.5 N force tolerance, with a documented
Unified Safety Level of at least 7.0 \cite{cfr312paiuni, pdac060s2030}. The
check is that the gate sign-off is a dated document, and the participant is
shown it at the baseline visit.

\textbf{Three. You are never the untested dose.} Daraxonrasib is escalated by a
standard 3+3 rule with staggered sentinel enrollment. A dose level opens only
after the level below has been cleared, and the second participant at a new
level is not dosed until the first participant's safety window has closed. The
check is the dose level printed on the participant's own assignment, and the
cohort record that shows what happened at the level below.

\textbf{Four. A human approves every motion.} The system is classified as a
Class II collaborative device under continuous human oversight. The on-premises
language model proposes; a deterministic safety gate tests the proposal; the
operating surgeon approves, modifies, or declines. No message travels from the
model to the arms. The check is the audit trail, in which every executed motion
carries a surgeon signature and no motion carries none.

\textbf{Five. Stopping is faster than harm.} The emergency-stop budget is at
most 3 ms across arms and at most 500 ms system-wide, on a 10 kHz heartbeat
broadcast bus \cite{ieee7009, iec8060127}. A human blink takes roughly 100 to
150 ms. Every arm is held before a blink completes. The check is that the
latency is measured on every procedure and reported, not assumed.

\textbf{Six. Your record cannot be edited.} Every model advisory, every gate
decision, every surgeon approval, and every safety-limit event is written to a
hash-chained audit trail bound to a deterministic seed and a specific repository
commit, under 21 CFR part 11 \cite{ecfr2026part11}. The check is that the
participant may request their own extract, and the chain either verifies or it
does not.

\textbf{Seven. Leaving costs you nothing.} Withdrawal is available at any visit,
without a reason, and standard-of-care oncology continues unchanged. The check
is in \S 8: four routes out are enumerated, and the data consequence of each is
stated, including the two consequences that are unwelcome.

\vspace{0.30em}

\begin{xltabular}{\textwidth}{L{0.7cm} L{4.6cm} Y}
\toprule
\textbf{\#} & \textbf{Commitment} & \textbf{The clause that makes it enforceable} \\
\midrule
\endfirsthead
\multicolumn{3}{@{}l}{\footnotesize\itshape Table continued from the previous page.}\\
\toprule
\textbf{\#} & \textbf{Commitment} & \textbf{The clause that makes it enforceable} \\
\midrule
\endhead
\midrule
\multicolumn{3}{r@{}}{\footnotesize\itshape Table continued on the next page.}\\
\endfoot
\bottomrule
\endlastfoot
1 & Nothing starts without you & Consent under 21 CFR \S 50.20 precedes every study procedure \\
2 & The robot is proven before you & Phase 0 gate: at least 1000 simulations across at least two independent frameworks, Unified Safety Level at least 7.0 \\
3 & You are never the untested dose & 3+3 escalation with sentinel staggering; a dose level opens only after the level below is cleared \\
4 & A human approves every motion & Class II collaborative device under continuous human oversight; no model-to-actuator path exists \\
5 & Stopping is faster than harm & Emergency stop at most 3 ms across arms and at most 500 ms system-wide \\
6 & Your record cannot be edited & Hash-chained audit trail under 21 CFR part 11, bound to a deterministic seed and a repository commit \\
7 & Leaving costs you nothing & Withdrawal at any visit without reason; standard care continues unchanged \\
\end{xltabular}

\vspace{0.30em}

\subsection{How to check a commitment, without taking anyone's word}

A commitment that cannot be checked is a slogan. Each of the seven above names
an artefact, and an artefact can be asked for. The list below is what to ask
for, and who holds it.

\vspace{0.30em}

\begin{xltabular}{\textwidth}{C{0.7cm} L{4.9cm} Y}
\toprule
\textbf{\#} & \textbf{Ask to see} & \textbf{Who holds it, and what a satisfactory answer looks like} \\
\midrule
\endfirsthead
\multicolumn{3}{@{}l}{\footnotesize\itshape Table continued from the previous page.}\\
\toprule
\textbf{\#} & \textbf{Ask to see} & \textbf{Who holds it, and what a satisfactory answer looks like} \\
\midrule
\endhead
\midrule
\multicolumn{3}{r@{}}{\footnotesize\itshape Table continued on the next page.}\\
\endfoot
\bottomrule
\endlastfoot
1 & The dated consent form, and the date of the first study-specific record & The site study coordinator. The consent date is on or before every study record date. A record that predates consent is a deviation and is reportable \\
2 & The Phase 0 gate sign-off sheet & The Sponsor-Investigator, through the site. It carries a date, the number of simulated procedures, the two framework names, and the Unified Safety Level as a number \\
3 & The dose level you are assigned to, and its clearance record & The site Principal Investigator. Your level is named in your consent; the level below it has a written clearance and closed sentinel windows \\
4 & The device classification and the oversight statement & The site. The device is Class II under continuous human oversight; the correct answer names the operating surgeon, not the software \\
5 & The most recent emergency-stop latency check for your room & The site biomedical engineering record. Two numbers, cross-arm and system-wide, each with a date and each under budget \\
6 & Your own audit extract & The Sponsor-Investigator, on request under \S 11.5. It arrives as a hash-chained record with a seed and a commit identifier, not as a summary \\
7 & The withdrawal paragraph of your consent form & The site. It states that standard care continues unchanged and that no reason is required \\
\end{xltabular}

\vspace{0.30em}

\noindent None of those requests is unusual and none of them requires a lawyer.
A site that cannot produce the artefact behind a commitment has not met the
commitment, whatever else it has done, and \S 11.5 names the escalation route
that applies in that case.

\subsection{What this paper is, and what it is not}

It is an independent research paper and a practical adoption guide, built
entirely from the author's repository sources. It is not medical advice, not
regulatory advice, and not endorsed by the FDA, the NIH, HHS, an Institutional
Review Board, the International Council for Harmonisation, or any sponsor,
contract research organisation, site, regulator, or medical society.

It is also not the informed consent form. A consent form has a required content
list, a required reading level, and a required signature block, and the National
Cancer Institute publishes a plain-language guide to what that document must
contain \cite{NCIConsent24}. This paper does none of that work. What it does is
prior: it sets out the twenty-one concerns people actually raise about surgical
robots, and it says, for each one, which clause of this specific protocol
answers it and how completely. A participant who has read this paper arrives at
the consent conversation with better questions.

Every figure in this paper derives from the author's repository sources and is
illustrative unless it is tied to a cited reference. Where a number is
simulated rather than measured in humans, \S 10 says so on the same line as the
number, and Figure 28 traces it to the source that produced it.

\subsection{How to read the thirty figures}

The paper carries thirty figures across five diagram types, and the type is not
decoration. Each type answers a different shape of question, and the choice of
type for a given figure is itself an argument about what kind of question is
being asked.

\vspace{0.30em}

\begin{xltabular}{\textwidth}{L{3.7cm} C{1.3cm} Y}
\toprule
\textbf{Diagram type} & \textbf{Count} & \textbf{The question it answers} \\
\midrule
\endfirsthead
\multicolumn{3}{@{}l}{\footnotesize\itshape Table continued from the previous page.}\\
\toprule
\textbf{Diagram type} & \textbf{Count} & \textbf{The question it answers} \\
\midrule
\endhead
\midrule
\multicolumn{3}{r@{}}{\footnotesize\itshape Table continued on the next page.}\\
\endfoot
\bottomrule
\endlastfoot
Mermaid-type & 9 & What happens, in what order, and who decides? Decisions in time \\
D2-type & 7 & How much of this is there, and how does it group? Containment and true grids \\
PlantUML-type & 5 & What exactly do you guarantee? Formal notation with a defined semantics \\
Graphviz-type & 5 & What depends on what, and can you prove it? Pure graph structure \\
Diagrams (Python)-type & 4 & Where does this physically live? Rooms, cabinets, cables \\
\end{xltabular}

\vspace{0.30em}

\noindent The monospace tag at the top left of every frame names the native
construct the figure reproduces, for example \texttt{flowchart TD} or
\texttt{digraph}, so a reader can carry any figure back to the platform it came
from. Every figure is vector line art drawn in the document itself. There is no
raster image anywhere in this paper.

\begin{pafloat}
\begin{pafig}[mermaid-type: flowchart TD]
\node[mmgoal,text width=34mm] (p)  at (0,0)      {\textbf{You}, the participant};
\node[mmin,text width=52mm]  (c1) at (0,-2.1)    {\textbf{1. Nothing starts without you}\\[1pt]{\scriptsize Consent signed before any study procedure, 21 CFR \S 50.20}};
\node[mmin,text width=52mm]  (c2) at (0,-4.5)    {\textbf{2. The robot is proven before you}\\[1pt]{\scriptsize Phase 0 gate: $\geq 1000$ simulations, $\geq 2$ frameworks, USL $\geq 7.0$}};
\node[mmin,text width=52mm]  (c3) at (0,-6.9)    {\textbf{3. You are never the untested dose}\\[1pt]{\scriptsize 3+3 escalation, sentinel stagger, DL1 cleared before DL2 opens}};
\node[mmin,text width=52mm]  (c4) at (0,-9.3)    {\textbf{4. A human approves every motion}\\[1pt]{\scriptsize Class II collaborative device, continuous oversight, no free running}};
\node[mmin,text width=52mm]  (c5) at (0,-11.7)   {\textbf{5. Stopping is faster than harm}\\[1pt]{\scriptsize $\leq 3$ ms cross-arm, $\leq 500$ ms system-wide guarantee}};
\node[mmin,text width=52mm]  (c6) at (0,-14.1)   {\textbf{6. Your record cannot be edited}\\[1pt]{\scriptsize Hash-chained audit trail, 21 CFR part 11, seed and commit bound}};
\node[mmin,text width=52mm]  (c7) at (0,-16.5)   {\textbf{7. Leaving costs you nothing}\\[1pt]{\scriptsize Withdraw at any time, standard care continues, \S 8 governs data}};
\node[mmdec,text width=15mm] (g1) at (6.6,-2.1)  {consent?};
\node[mmdec,text width=15mm] (g2) at (6.6,-4.5)  {gate?};
\node[mmdec,text width=15mm] (g3) at (6.6,-6.9)  {cohort?};
\node[mmdark,text width=28mm](st) at (12.4,-4.5) {No study activity occurs at all};
\node[mmstep,text width=52mm](d30)at (0,-18.9)   {Day 30, primary safety window\\Clavien-Dindo grading complete};
\node[mmgoal,text width=52mm](d90)at (0,-21.1)   {Day 90, R0 pathology and 90-day mortality reported to you directly};
\draw[mmedge] (p)  -- (c1);
\draw[mmedge] (c1) -- (c2);
\draw[mmedge] (c2) -- (c3);
\draw[mmedge] (c3) -- (c4);
\draw[mmedge] (c4) -- (c5);
\draw[mmedge] (c5) -- (c6);
\draw[mmedge] (c6) -- (c7);
\draw[mmedge] (c7) -- (d30);
\draw[mmedge] (d30) -- (d90);
\draw[mmedgeb] (c1) -- node[mmlabel,above]{\scriptsize consent on file} (g1);
\draw[mmedgeb] (c2) -- node[mmlabel,above]{\scriptsize gate signed} (g2);
\draw[mmedgeb] (c3) -- node[mmlabel,above]{\scriptsize level cleared} (g3);
\draw[mmedged] (g1) to[out=0,in=150,looseness=0.9] node[mmlabel,above,pos=0.55]{no} (st.north west);
\draw[mmedged] (g2) -- node[mmlabel,above]{no} (st);
\draw[mmedged] (g3) to[out=0,in=210,looseness=0.9] node[mmlabel,below,pos=0.55]{no} (st.south west);
\node[font=\scriptsize\sffamily\itshape,text=protogray,anchor=north west,align=left]
  at (-3.0,-22.6) {Each gate closes behind the commitment above it. A ``no'' at any gate stops every downstream activity,\\which is what makes the commitment enforceable rather than aspirational.};
\end{pafig}
\vspace{-0.7cm}
\figcaption{Figure 1. The seven commitments made to the participant, from the first conversation\\
through consent, the Phase 0 gate, and dose assignment, to the day-30 safety\\
window and the day-90 pathology result that is reported back to them directly.}
\end{pafloat}

\subsection{Accountability, stated once, at the front}

The single most common accountability question in the surveyed literature is
the plainest one: if something goes wrong, who is to blame
\cite{Jauniaux2025, BrarRAS2024X, WuSemiAI2026}. When a human surgeon errs, the
chain of liability is familiar. When a semi-autonomous movement or a
model-generated suggestion contributes to harm, participants report being told
that the surgeon, the hospital, the software developer, and the device
manufacturer are all involved, which is not an answer.

This protocol answers it with a rule rather than a list: exactly one accountable
party per failure class, named in advance. Section 11 sets out the full
responsibility matrix, and Figure 29 renders it with a final column giving the
person the participant should actually call. The rule is checkable, because a
matrix with two accountable parties in one row would fail it, and the matrix is
published.

Figure 2 states the same structure as a graph. The chart is rooted at the
participant rather than at the sponsor, and it is acyclic, which means that no
body in it reviews itself. Three edges point back at the root, and only three,
because three is what the protocol actually owes the participant as an
affirmative obligation to inform.

\begin{pafloat}
\begin{pafig}[graphviz-type: digraph, rooted DAG]
\node[gvkey,text width=24mm] (you) at (0,0) {\textbf{YOU}\\the participant};
\node[gvsoft,text width=22mm] (sur) at (-6.9,-3.4) {Operating surgeon\\approves every motion};
\node[gvnode,text width=22mm] (ana) at (-2.3,-3.4) {Anaesthetist};
\node[gvsoft,text width=22mm] (ism) at (2.3,-3.4) {Independent safety monitor, may stop};
\node[gvnode,text width=22mm] (scr) at (6.9,-3.4) {Scrub and circulating team};
\node[gvsoft,text width=22mm] (spi) at (-2.5,-7.0) {Site principal investigator};
\node[gvmid,text width=22mm]  (spo) at (2.5,-7.0)  {Sponsor-Investigator ChemicalQDevice};
\node[gvgray2,text width=22mm](dsm) at (-5.0,-10.6) {DSMB, may halt the study};
\node[gvgray2,text width=22mm](pai) at (0,-10.6) {Physical AI Safety Review Committee};
\node[gvgray2,text width=22mm](irb) at (5.0,-10.6)  {IRB, you may write directly};
\node[gvbox,fill=pagrayd,text width=26mm] (fda) at (0,-13.4) {FDA, IND and IDE};
\node[gvboxd,text width=54mm] (duty) at (0,-15.4) {The obligation every path above is discharging: your safety first, then the evidence};
\begin{scope}[on background layer]
\node[gvcluster,fit=(sur)(ana)(ism)(scr)] (Croom) {};
\node[gvcluster,fit=(spi)(spo)] (Cstudy) {};
\node[gvcluster2,fit=(dsm)(pai)(irb)] (Cind) {};
\end{scope}
\node[gvctitle,anchor=south west] at (Croom.north west) {in the room on the day};
\node[gvctitle,anchor=south west] at (Cstudy.north west) {accountable for the study};
\node[gvctitle,anchor=south west] at (Cind.north west) {independent of the sponsor};
\draw[gvedge] (you) to[out=-150,in=75,looseness=0.9] node[mmlabel,pos=0.62]{consents to} (sur);
\draw[gvedge] (you) to[out=-60,in=105,looseness=0.9] (ism);
\draw[gvedge] (sur) to[out=-75,in=160,looseness=0.9] node[mmlabel,pos=0.55]{reports} (spi);
\draw[gvedge] (ana) -- (spi);
\draw[gvedge] (scr) to[out=-105,in=20,looseness=0.9] (spi);
\draw[gvedge] (ism) to[out=-105,in=30,looseness=0.9] node[mmlabel,pos=0.5]{escalates} (spi);
\draw[gvedged] (ism) to[out=-75,in=80,looseness=0.9] node[mmlabel,pos=0.55]{may escalate past the site} (pai);
\draw[gvedge] (spi) -- node[mmlabel,above]{reports} (spo);
\draw[gvedge] (spo) to[out=-135,in=45,looseness=0.9] (dsm);
\draw[gvedge] (spo) to[out=-105,in=60,looseness=0.9] (pai);
\draw[gvedge] (spo) to[out=-60,in=105,looseness=0.9] (irb);
\draw[gvedged] (dsm) to[out=-60,in=170,looseness=0.9] (fda);
\draw[gvedged] (pai) -- (fda);
\draw[gvedged] (irb) to[out=-120,in=10,looseness=0.9] (fda);
\draw[gvedgeb] (fda) -- (duty);
% The three affirmative obligations returning to the root, routed in clear
% channels outside every cluster so no edge crosses a node or a cluster border.
\draw[gvedgeb] (spo.east) -- (11.4,-7.0) -- (11.4,0)
  node[mmlabel,midway,align=left,anchor=west,xshift=2pt]{must tell you:\\version change,\\new risk, and\\your own results} -- (you.east);
\draw[gvedgeb] (dsm.west) -- (-11.4,-10.6) -- (-11.4,-0.6)
  node[mmlabel,pos=0.62,align=right,anchor=east,xshift=-2pt]{must tell you:\\a study pause} -- (you.west);
\draw[gvedgeb] (spi.west) -- (-8.9,-7.0) -- (-8.9,-1.1)
  node[mmlabel,pos=0.55,align=right,anchor=east,xshift=-2pt]{must tell you:\\any adverse event\\affecting you} -| (you.south);
\end{pafig}
\vspace{-0.7cm}
\figcaption{Figure 2. The accountability chain drawn as a rooted directed acyclic graph with\\
the participant at the root, showing the people in the room, the bodies accountable\\
for the study, the independent reviewers, and the three obligations that point back.}
\end{pafloat}

\noindent The graph is acyclic by construction, and the property is load-bearing
rather than aesthetic: a cycle would mean that some body in the chart reviews
its own work, which is precisely the failure the accountability concern
describes. The three blue edges returning to the root are the affirmative
obligations to inform. They are not courtesies, and \S 11.5 names the escalation
route that applies if one of them is not met.

\subsection{What this paper does not promise}

Advocacy that promises too much is worth less than advocacy that draws its own
boundary, so the boundary is drawn here rather than buried in \S 10.

This paper does not promise that the operation will succeed. It is a Phase 1
study, and the primary questions are feasibility and safety, not cure. A
participant who joins expecting a better oncologic outcome than a conventional
resection is joining for a reason the design cannot support, and \S 4.1 states
the reason the design cannot support it.

It does not promise that the platform is better than a human-only Whipple. No
comparison of that kind is possible at this sample size, and \S 10.5 gives the
arithmetic that shows why: with at most 18 participants and no randomised
control, the study can detect a gross safety signal and nothing finer. Where
this paper cites a comparator number, the comparator is a published series and
is labelled C in the provenance table, never a result of this study.

It does not promise that nothing will hurt. A Whipple is one of the largest
operations in general surgery under any technique. The relevant claim is
narrower and is the one the protocol can keep: that the additional risk
introduced by the platform, over and above the operation itself, is bounded by
limits that are enforced in hardware and measured on every case.

It does not promise that every concern is answered. Five of the twenty-one are
answered by governance or disclosure alone, and \S 3.7 marks each of those five
as such rather than dressing it as a guarantee. Two cost items in \S 12.3 remain
unsettled at the time of writing, and they are printed unsettled.

Finally, it does not promise that the participant's judgement should defer to
the protocol's. The purpose of the document is the opposite: to put the
participant in a position to disagree with it in specific terms.
